\SourceLeaf{10}
\begin{center}
{\Large\textsc{PRAEFATIO:}}

\medskip

\BookRule{0.07\linewidth}
\end{center}

\noindent Disquisitiones in hoc opere contentae ad eam Matheseos partem pertinent, quae circa numeros integros versatur, fractis plerumque, surdis semper exclusis. Analysis indeterminata quam vocant seu Diophantaea, quae ex infinitis solutionibus problemati indeterminato satisfacientibus eas seligere docet, quae per numeros integros aut saltem rationales absolvuntur (plerumque ea quoque conditione adiecta ut sint positivi) non est illa disciplina ipsa, sed potius pars eius valde specialis, ad eamque ita fere se habet, ut ars aequationes reducendi et solvendi (Algebra) ad universam Analysin. Nimirum quemadmodum ad \emph{Analyseos} ditionem referuntur omnes quae circa quantitatum affectiones generales institui possunt disquisitiones: ita numeri integri (fractique quatenus per integros determinantur) obiectum proprium \textsc{ARITHMETICAE} constituunt. Sed quum ea, quae Arithmetices nomine vulgo traduntur, vix ultra artem numerandi et calculandi (i.~e. numeros per signa idonea e.~g. secundum systema decadicum exhibendi, operationesque arithmeticas perficiendi) extendantur, adiectis nonnullis quae vel ad Arithmeticam omnino non pertinent (ut doctrina de logarithmis) vel saltem numeris integris non sunt propria sed ad omnes quantitates patent: e re esse videtur, duas Arithmeticae partes distinguere, illaque ad Arithmeticam elementarem referre, omnes autem disquisitiones generales de numerorum integrorum affectionibus propriis \emph{Arithmeticae Sublimiori}, de qua sola hic sermo erit. vindicare.

Pertinent ad Arithmeticam Sublimiorem ea, quae Euclides in Elementis L. VII sqq. elegantia et rigore apud veteres consuetis tradidit: attamen ad prima initia huius scientiae limitantur. Diophanti opus celebre, quod totum problematis \SourceLeaf{11}indeterminatis dicatum est, multas quaestiones continet, quae propter difficultatem suam artificiorumque subtilitatem de auctoris ingenio et acumine existimationem haud mediocrem suscitant, praesertim si subsidiorum quibus illi uti licuit tenuitatem consideres. At quum haec problemata dexteritatem quandam potius scitamque tractationem, quam principia profundiora postulent, praetereaque nimis specialia sint raroque ad conclusiones generaliores deducant: hic liber ideo magis epocham in historia Matheseos constituere videtur, quod prima artis characteristicae et Algebrae vestigia sistit, quam quod Arithmeticam Sublimiorem inventis novis auxerit. Longe plurima recentioribus debentur, inter quos pauci quidem sed immortalis gloriae viri P. \textsc{de Fermat}, L. \textsc{Euler}, L. \textsc{La Grange}, A. M. \textsc{Le Gendre} (ut paucos alios praeteream) introitum ad penetralia huius divinae scientiae aperuerunt, quantisque divitiis abundent patefecerunt. Quaenam vero inventa a singulis his geometris profecta sint, hic enarrare supersedeo, quum e praefationibus Additamentorum quibus ill. La Grange Euleri Algebram ditavit operisque mox memorandi ab ill. Le Gendre nuper editi cognosci possint, insuperque pleraque locis suis in his Disquisitionibus Arithmeticis laudentur.

Propositum huius operis, ad quod edendum iam annos abhinc quinque publice fidem dederam, id fuit, ut disquisitiones ex Arithmetica Sublimiori, quas partim ante id tempus partim postea institui, divulgarem. Ne quis vero miretur, scientiam hic a primis propemodum initiis repetitam, multasque disquisitiones hic denuo resumtas esse, quibus alii operam suam iam navarunt. monendum esse duxi, me, quum primum initio a. 1795 huic disquisitionum generi animum applicavi, omnium quae quidem a recentioribus in hac arena elaborata fuerint ignarum, omniumque subsidiorum per quae de his quidpiam comperire potuissem expertem fuisse. Scilicet in alio forte labore tunc occupatus. casu incidi in eximiam quandam veritatem arithmeticam (fuit autem ni fallor theorema art. 108), quam quum et per se pulcherrimam aestimarem et cum maioribus connexam esse suspicarer. summa qua potui contentione in id incubui, ut principia quibus inniteretur perspicerem, demonstrationemque rigorosam nanciscerer. Quod postquam tandem ex voto successisset. illecebris harum quaestionum ita fui implicatus, ut eas deserere non potuerim; quo pacto, dum alia semper ad alia viam sternebant, ea quae in quatuor primis Sectionibus huius operis traduntur ad maximam partem absoluta erant, antequam de aliorum geometrarum laboribus similibus quidquam vi\SourceLeaf{12}dissem. Dein copia mihi facta, horum summorum ingeniorum scripta evolvendi, maiorem quidem partem meditationum mearum rebus dudum transactis impensam esse agnovi: sed eo alacrior, illorum vestigiis insistens, Arithmeticam ulterius excolere studui; ita variae disquisitiones institutae sunt, quarum partem Sectiones V, VI et VII tradunt. Postquam interiecto tempore consilium de fructibus vigiliarum in publicum edendis cepi: eo lubentius, quod plures optabant, mihi persuaderi passus sum, ne quid vel ex illis investigationibus prioribus supprimerem, quod tum temporis liber non habebatur, ex quo aliorum geometrarum labores de his rebus, in Academiarum Commentariis sparsi, edisci potuissent; quod multae ex illis omnino novae et pleraeque per methodos novas tractatae erant; denique quod omnes tum inter se tum cum disquisitionibus posterioribus tam \emph{arcto} nexu cohaerebant, ut ne nova quidem satis commode explicari possent, nisi reliquis ab initio repetitis.

Prodiit interea opus egregium viri iam antea de Arithmetica Sublimiori magnopere meriti, \emph{Le Gendre Essai d'une théorie des nombres, Paris a. VI,} in quo non modo omnia quae hactenus in hac scientia elaborata sunt diligenter collegit et in ordinem redegit, sed permulta insuper nova de suo adiecit. Quum hic liber serius ad manum mihi pervenerit, postquam maxima operis pars typis iam exscripta esset; nullibi, ubi rerum analogia occasionem dare potuisset, eius mentionem iniicere licuit; de paucis tantummodo locis quaedam observationes in Additamentis adiungere necessarium videbatur, quas vir humanissimus et candidissimus benigne ut spero interpretabitur.

Inter impressionem huius operis, quae pluries interrupta variisque impedimentis usque in quartum annum protracta est, non modo eas investigationes, quas quidem iam antea susceperam, sed quarum promulgationem in aliud tempus differre constitueram, ne liber nimis magnus evaderet, ulterius continuavi, sed plures etiam alias novas aggressus sum. Plures quoque, quas ex eadem ratione leviter tantum attigi, quum tractatio uberior minus necessaria videretur (e. g. eae quae in artt. 37, 82 sqq. aliisque locis traduntur), postea resumtae sunt, disquisitionibusque generalioribus quae luce perdignae videntur locum dederunt (Conf. etiam quae in Additamentis de art. 306 dicuntur). Denique quum liber praesertim propter amplitudinem Sect. V in longe maius quam exspectaveram volumen excres\SourceLeaf{13}ceret, plura quae ab initio ei destinata erant, interque ea totam Sectionem \emph{octavam} (quae passim iam in hoc volumine commemoratur, atque tractationem generalem de congruentiis algebraicis cuiusvis gradus continet) resecare oportuit. Haec omnia. quae volumen huic aequale facile explebunt, publici iuris fient, quamprimum occasio aderit.

Quod. in pluribus quaestionibus difficilibus, demonstrationibus syntheticis usus sum. analysinque per quam erutae sunt suppressi, imprimis brevitatis studio tribuendum est. cui quantum fieri poterat consulere oportebat.

Theoria divisionis circuli, sive polygonorum regularium, quae in Sect. VII tractatur, \emph{ipsa} quidem \emph{per se} ad Arithmeticam non pertinet, attamen eius \emph{principia} unice ex Arithmetica Sublimiori petenda sunt: quod forsan geometris tam inexspectatum erit, quantum veritates novas, quas ex hoc fonte haurire licuit, ipsis gratas fore spero.

Haec sunt. de quibus lectorem praemonere volui. De rebus ipsis non meum est iudicare. Nihil equidem magis opto, quam ut iis, quibus scientiarum incrementa cordi sunt, placeant, quae vel hactenus desiderata explent, vel aditum ad nova aperiunt.

\BookRule{0.16\linewidth}
